%% parts/p05/diagonalization.tex --- Part 5: Diagonalization and the First
%% Incompleteness Theorem (Phase-3 pilot migration of v1.1
%% ch:first-incompleteness; load-bearing apparatus of v1.1
%% ch:arithmetic-coding restated in Section 1).  Body = statements only;
%% proofs live in annex/p05.tex.
\chapter{Diagonalization and the First Incompleteness Theorem}
\label{sec:p05-diagonalization}
\chaptermark{Diagonalization and Incompleteness}

This part separates three levels throughout: arithmetic in the standard
structure, formal derivability in a theory \(T\), and metalogical reasoning
about \(T\).  G\"odel's original paper is \cite{Godel1931}; modern
treatments include \cite{Shoenfield1967,BoolosBurgessJeffrey2007}.

\section{Arithmetical prerequisites}
\label{sec:p05-diagonalization-prereq}

The entries of this section restate, from the v1.1 arithmetic chapter,
exactly the apparatus that the incompleteness results load as hypotheses:
the language and numerals, the fixed calculus, \(\mathsf Q\) and \(\PA\),
the registered order abbreviation with its two comparison lemmas, the
representability vocabulary, and the fixed proof predicate.

\begin{hbdefinition}{def:p05-language-arithmetic}{Arithmetic language and standard structure}
The \emph{language of arithmetic} is \(\Lang_A=\{0,S,+,\times\}\), where
\(0\) is a constant, \(S\) is unary, and \(+\) and \(\times\) are binary
function symbols.  The \emph{standard structure of arithmetic} is
\(\mathcal N_{\mathrm{ar}}=(\N,0,S,+,\times)\) with the usual operations.  For
\(n\in\N\), the \emph{numeral} for \(n\) is the closed term
\(\bar n=S^n0\); for a tuple \(\vec a=(a_1,\ldots,a_k)\),
\(\overline{\vec a}\) is the tuple of numerals
\((\bar a_1,\ldots,\bar a_k)\).
\end{hbdefinition}

\begin{hbconvention}{def:p05-turnstile}{The fixed calculus and its turnstile}
Throughout this part, \(T\derives\sigma\) says that a finite derivation of
\(\sigma\) from the axioms of \(T\) exists in the fixed classical explosive
first-order calculus with equality recorded as sheet \texttt{nd-fol} in the
calculus registry; under this declared local abbreviation the per-system
subscript of the registered turnstile is dropped.  \(T\notderives\sigma\)
denies it.  Soundness and semantic completeness of the fixed calculus are
as recorded by the classical adequacy entries of the first-order part.
\end{hbconvention}

\AddToHookNext{env/definition/end}{\label{def:p05-robinson-q}}
\begin{inheritedrestorationscope}{SRC-09099-09111}
\begin{definition}[Robinson arithmetic]
Robinson arithmetic, written \(\mathsf Q\), is the finitely axiomatized
\(\Lang_A\)-theory consisting of the universal closures of:
\[
\begin{array}{ll}
  Sx\neq0, & Sx=Sy\to x=y,\\
  x\neq0\to\exists y\,x=Sy, & x+0=x,\\
  x+Sy=S(x+y), & x\times0=0,\\
  x\times Sy=(x\times y)+x. &
\end{array}
\]
It has no induction scheme.
\end{definition}

\end{inheritedrestorationscope}


\begin{hbdefinition}{def:p05-peano-arithmetic}{Peano arithmetic}
Peano arithmetic, written \(\PA\), is \(\mathsf Q\) together with every
instance of the induction scheme
\[
  \bigl(\varphi(0,\vec z)\land
  \forall x(\varphi(x,\vec z)\to\varphi(Sx,\vec z))\bigr)
  \to\forall x\,\varphi(x,\vec z),
\]
where \(\varphi(x,\vec z)\) is any \(\Lang_A\)-formula and the parameters
\(\vec z\) remain universally quantified in the closure.
\end{hbdefinition}

\begin{hbconvention}{def:p05-leq}{The registered order abbreviation}
For arithmetic terms \(r,t\), the registered object-language abbreviation
\(\leq\) expands as \(r\leq t\defeq\exists z\,(z+r=t)\), where \(z\) is the
least variable absent from \(r,t\); for \(x\notin\operatorname{Var}(t)\),
the marked bounded quantifiers
\(\exists x\leq t\,\varphi\) and \(\forall x\leq t\,\varphi\) abbreviate
\(\exists x\,(x\leq t\land\varphi)\) and
\(\forall x\,(x\leq t\to\varphi)\) and are counted as bounded.  The
witness \(z\) sits to the \emph{left} of \(+\) in \(z+r=t\): this choice
is what lets the two comparison lemmas below rest only on the successor
axioms and the right-recursion axioms \(x+0=x\) and \(x+Sy=S(x+y)\) for
addition.
\end{hbconvention}

The incompleteness argument of Rosser evaluates bounded numerical
comparisons inside \(\mathsf Q\).  The two lemmas below supply exactly
those evaluations for the registered abbreviation \(\leq\).

\begin{hblemma}{thm:p05-q-bounded-cases}{proved=annex:p5.1}{Bounded-case enumeration in \(\mathsf Q\)}
For every \(n\in\N\),
\[
  \mathsf Q\derives\forall y\,\bigl(y\leq\bar n\to
  (y=\bar0\lor y=\bar1\lor\dots\lor y=\bar n)\bigr).
\]
\end{hblemma}

\begin{hblemma}{thm:p05-q-dichotomy}{proved=annex:p5.2}{Dichotomy in \(\mathsf Q\)}
For every \(n\in\N\),
\[
  \mathsf Q\derives\forall y\,
  \bigl(y\leq\bar n\lor\bar n\leq y\bigr).
\]
\end{hblemma}

\begin{hbdefinition}{def:p05-numeralwise-representation}{Numeralwise representation of a relation}
Let \(T\supseteq\mathsf Q\), and let \(R\subseteq\N^k\).  A formula
\(\varphi_R(x_1,\ldots,x_k)\) \emph{numeralwise represents} \(R\) in \(T\)
if, for every \(\vec n\in\N^k\): if \(R(\vec n)\) then
\(T\derives\varphi_R(\overline{\vec n})\), and if \(\neg R(\vec n)\) then
\(T\derives\neg\varphi_R(\overline{\vec n})\).  The two clauses concern each
standard tuple of numerals; they do not assert that \(T\) proves a
universally quantified decision statement.
\end{hbdefinition}

\begin{hbdefinition}{def:p05-function-representation}{Representation of a function}
Let \(T\supseteq\mathsf Q\).  A formula \(\gamma_f(\vec x,y)\)
\emph{represents} a total function \(f:\N^k\to\N\) in \(T\) if, for every
\(\vec n\in\N^k\),
\[
  T\derives\forall y\,
  \bigl(\gamma_f(\overline{\vec n},y)
  \leftrightarrow y=\overline{f(\vec n)}\bigr).
\]
\end{hbdefinition}

\begin{hbtheorem}{thm:p05-representability}{cited=Shoenfield1967}{Representability theorem}
Every primitive recursive relation is numeralwise representable in
\(\mathsf Q\) by a \(\Sigma^0_1\) formula, and every primitive recursive
function is representable in \(\mathsf Q\).  Every computably enumerable
relation is definable in \(\mathcal N_{\mathrm{ar}}\) by a \(\Sigma^0_1\) formula.
\end{hbtheorem}

The construction represents the initial functions directly, preserves
representability through composition and primitive recursion by
quantifying over coded finite computation sequences, and supplies a
finite accepting computation as the existential witness in the computably
enumerable case; \cite{Shoenfield1967} carries the complete coding proof.

\begin{hbdefinition}{def:p05-goedel-numbering}{Godel numbering@G\"odel numbering}
A \emph{G\"odel numbering} for the chosen formal language and calculus is
an injective effective map from symbols, finite strings, formulas, and
finite derivations to natural numbers such that the basic syntactic
operations and predicates required below are primitive recursive on
codes.
\end{hbdefinition}

The exact numerical values are conventional; the effectiveness properties
are substantive.  Fix one such numbering.  Write \(\ulcorner E\urcorner\)
for the code of the expression or derivation \(E\); in an arithmetic
formula, \(\overline{\ulcorner E\urcorner}\) is the numeral denoting that
code.

\begin{hbdefinition}{def:p05-ce-theory}{Computably enumerable theory}
Relative to the fixed G\"odel numbering, a theory
\(T\subseteq\mathrm{Sent}(\Lang_A)\) is \emph{computably enumerable} iff
\(\{\ulcorner\sigma\urcorner:\sigma\in T\}\) is a computably enumerable
set of natural numbers.  This term classifies the given set of axioms
\(T\); it does not mean that \(T\) is deductively closed, and it does not
define \(T\) as the set of its derivable consequences.  The distinction
is required because the proof predicate below verifies uses of members of
\(T\) by finite enumeration-stage certificates.
\end{hbdefinition}

\begin{hbdefinition}{def:p05-proof-predicate}{Proof and provability predicates}
Let \(T\) be computably enumerable and contain \(\mathsf Q\).  Fix a
primitive recursive relation \(\mathrm{Proof}_T(p,x)\) meaning that
\(p\) codes a correct \(T\)-proof whose concluding sentence has code
\(x\), with enumeration-stage certificates verifying the uses of axioms
of \(T\).  By the \(\Sigma^0_1\) clause of
Theorem~\ref{thm:p05-representability}, fix the same symbol for a
specific \(\Sigma^0_1\) formula that numeralwise represents this relation
in \(\mathsf Q\).  Define
\[
  \mathrm{Prov}_T(x)\defeq\exists p\,\mathrm{Proof}_T(p,x).
\]
The additional existential quantifier merges with the initial existential
block of the chosen representation, so \(\mathrm{Prov}_T\) is a
\(\Sigma^0_1\) formula by the syntactic hierarchy definition.  A single
primitive recursive sequence coding (G\"odel's \(\beta\) function) stands
behind the \(\Sigma^0_1\) clause, serving machine tables, computation
histories, and proof sequences alike.
\end{hbdefinition}

\begin{hbwarning}{def:p05-mention-warning}{Mentioning a sentence inside arithmetic}
The expression \(\mathrm{Prov}_T(\varphi)\) is ill-typed if \(\varphi\)
is a formula rather than a number term.  The correct expression is
\(\mathrm{Prov}_T(\overline{\ulcorner\varphi\urcorner})\), abbreviated
only after the coding convention has been declared
(Section~\ref{sec:p05-provability-truth-01}).
\end{hbwarning}

\paragraph{Open inherited content.}
The exact source presentations for the nonstandard-model entries,
provable \(\Sigma^0_1\)-completeness, the arithmetization-of-syntax proof
sketch, the worked addition and sequence-code illustrations, and all three
exercise prompts remain owed in their original source order, standing, and
full content.  Chapter~\ref{sec:p05-models-and-interpretability},
Definition~\ref{def:p04-numerical-coding},
Lemma~\ref{thm:p04-goedel-beta},
Theorem~\ref{thm:p04-bounded-verification}, and the entries above are useful
cross-references but do not replace, consolidate, summarize, or discharge
any of those inherited presentations.  Their destinations and separate
source, static, build, visual, and human-content verification remain open.

\section{The diagonal lemma}
\label{sec:p05-diagonalization-01}

Let \(\mathrm{diag}(n)\) be the code of the result of substituting the
numeral \(\bar n\) for the distinguished free variable in the
one-variable formula with code \(n\), with an arbitrary default value on
nonformula codes.  The function \(\mathrm{diag}\) is primitive recursive.

\begin{hbtheorem}{thm:p05-diagonal-lemma}{cited=BoolosBurgessJeffrey2007}{Diagonal lemma}
Let \(T\) be any theory extending \(\mathsf Q\).  For every
\(\Lang_A\)-formula \(\psi(x)\) with at most the displayed free
variable, there is a sentence \(G\) such that
\[
  T\derives G\leftrightarrow
  \psi(\overline{\ulcorner G\urcorner}).
\]
\end{hbtheorem}

See \cite{BoolosBurgessJeffrey2007} for a complete proof, including the
primitive-recursive verification and numeralwise representability of
\(\mathrm{diag}\); the fixed-point construction behind the lemma is
retained as commentary in
\hyperref[annex:p5-diagonal-commentary]{the Part-5 annex}.

\begin{hbmethod}{def:p05-diagonal-method}{The diagonal method}
To obtain a sentence provably equivalent to a property of its own code:
fix a formula \(D(x,y)\) representing the substitution function
\(\mathrm{diag}\) in \(\mathsf Q\); given the target property
\(\psi(x)\), form \(\beta(x)\defeq\exists y\,(D(x,y)\land\psi(y))\);
and evaluate \(\beta\) at the numeral for its own code, setting
\(G\defeq\beta(\bar b)\) with \(b=\ulcorner\beta\urcorner\).
Representability of \(\mathrm{diag}\) then delivers
\(\mathsf Q\derives G\leftrightarrow
\psi(\overline{\ulcorner G\urcorner})\).  The construction is spelled
out in \hyperref[annex:p5-diagonal-commentary]{the Part-5 annex}.
\end{hbmethod}

\begin{hbwarning}{def:p05-selfreference}{Self-reference without self-inscription}
The diagonal lemma does not create a sentence that literally contains its
own inscription.  It creates a sentence provably equivalent to the result
of applying \(\psi\) to the numeral for that sentence's code.
\end{hbwarning}

\section{Consistency and soundness conditions}
\label{sec:p05-diagonalization-02}

\begin{hbdefinition}{def:p05-consistency}{Consistency}
A theory \(T\) is \emph{consistent} iff there is no sentence \(\sigma\)
such that both \(T\derives\sigma\) and \(T\derives\neg\sigma\).  In the
fixed classical explosive calculus, this is equivalent to
\(T\notderives\bottom\).  If \(T\derives 0\ne S0\), then it is also
equivalent to \(T\notderives 0=S0\); in particular this specialization
applies when \(T\supseteq\mathsf Q\).  Without the displayed
derivability hypothesis, nonderivability of \(0=S0\) is not a general
test for consistency.
\end{hbdefinition}

A theory \(T\) is \emph{deductively complete} when, for every sentence
\(\sigma\) in its language, either \(T\derives\sigma\) or
\(T\derives\neg\sigma\).  This property of a theory is distinct from
semantic completeness of a calculus, while soundness and completeness of
the fixed calculus identify it with the corresponding quantified semantic
condition.

\begin{hbdefinition}{def:p05-sigma-soundness}{Sigma-0-1-soundness@\(\Sigma^0_1\)-soundness}
A theory \(T\) is \emph{\(\Sigma^0_1\)-sound} iff, whenever \(\sigma\) is
a \(\Sigma^0_1\) sentence and \(T\derives\sigma\), then
\(\mathcal N_{\mathrm{ar}}\entails\sigma\).
\end{hbdefinition}

\begin{hbdefinition}{def:p05-omega-consistency}{Omega-consistency}
A theory \(T\) is \emph{omega-consistent} iff there is no formula
\(\varphi\) with \(\FV(\varphi)\subseteq\{x\}\) such that
\[
  T\derives\exists x\,\varphi(x)
  \quad\text{and}\quad
  T\derives\neg\varphi(\bar n)
  \text{ for every }n\in\N.
\]
\end{hbdefinition}

Both omega-consistency and \(\Sigma^0_1\)-soundness imply consistency,
but consistency alone is weaker.  Rosser's theorem removes the stronger
hypothesis from the incompleteness conclusion.

\section{The G\"odel sentence}
\label{sec:p05-diagonalization-03}

Let \(T\) be a computably enumerable theory extending \(\mathsf Q\).
Apply Theorem~\ref{thm:p05-diagonal-lemma}, in the complexity-preserving
form discussed below display (G), to
\(\psi(x)=\neg\mathrm{Prov}_T(x)\).  Obtain a \(\Pi^0_1\) sentence
\(G_T\) satisfying
\[
  T\derives G_T\leftrightarrow
  \neg\mathrm{Prov}_T(\overline{\ulcorner G_T\urcorner}).
  \tag{G}
\]
The fixed-point biconditional is in fact derivable already in
\(\mathsf Q\).  The complexity-preservation clause is a standard
refinement of the coding construction: choose a suitable representation
of substitution on codes and put the resulting fixed point back into
prenex form.  The fixed point of a \(\Pi^0_1\) formula may then be chosen
\(\Pi^0_1\).  This clause is cited rather than proved here; see
\cite{Shoenfield1967}.

\begin{hbtheorem}{thm:p05-goedel-first}{sketched=annex:p5.3}{Godel's first incompleteness theorem@G\"odel's first incompleteness theorem, original form}
Let \(T\) be a computably enumerable theory extending \(\mathsf Q\).
\begin{enumerate}
  \item If \(T\) is consistent, then \(T\notderives G_T\).
  \item If \(T\) is \(\Sigma^0_1\)-sound, then
        \(T\notderives\neg G_T\).
  \item Omega-consistency also suffices for
        \(T\notderives\neg G_T\).
\end{enumerate}
Consequently, if \(T\) is omega-consistent or \(\Sigma^0_1\)-sound, it is
not deductively complete.
\end{hbtheorem}

\begin{hbwarning}{def:p05-first-theorem-scope}{What the first theorem does not say}
The theorem does not say that every formal theory is deductively
incomplete, that \(G_T\) is unprovable in every stronger theory, or that
no deductively incomplete theory can prove any consistency statement.
The hypotheses and the subscript \(T\) are essential.
\end{hbwarning}

\section{Rosser's strengthening}
\label{sec:p05-diagonalization-04}

\begin{hbtheorem}{thm:p05-rosser}{sketched=annex:p5.4}{Rosser's theorem}
Let \(T\) be a consistent computably enumerable theory extending
\(\mathsf Q\).  Then there is a sentence \(R_T\) such that
\[
  T\notderives R_T
  \qquad\text{and}\qquad
  T\notderives\neg R_T.
\]
Thus consistency alone suffices for deductive incompleteness.
\end{hbtheorem}

The original source is \cite{Rosser1936}.  The Rosser sentence compares
proof codes with the registered bounded order
(Convention~\ref{def:p05-leq}); the comparison lemmas of
Section~\ref{sec:p05-diagonalization-prereq} carry those steps inside
\(\mathsf Q\).

\begin{hbcorollary}{thm:p05-theoremhood-undecidable}{sketched=annex:p5.5}{Undecidability of theoremhood}
If \(T\) is a consistent computably enumerable extension of
\(\mathsf Q\), then the set of codes of theorems of \(T\) is computably
enumerable but undecidable.
\end{hbcorollary}

An alternative proof reduces the self-application domain set to
theoremhood by uniform representability.  The completion route taken in
the annex makes explicit why independence of one fixed sentence does not
by itself contradict decidability of the theorem set.
